%% comparison.tex
%% Rewritten "Comparison to previous Pop III simulations" section.
%% Replace \section{Comparison...} through the end of that section
%% in oja_template.tex with this content.
%% ============================================================

\section{Comparison to previous Pop~III simulations}
\label{sec:disc_comparison}

Our simulation combines three features that, taken together, are
uncommon in the existing literature: fully cosmological initial
conditions, radiation-MHD with resolved individual protostellar
formation (STARFORGE), and $15\kyr$ of post-collapse evolution in a
halo progenitor of a Milky Way-mass system.
Below we compare our results in each dimension to prior work.

% -----------------------------------------------------------
\subsection*{Fragmentation and sink counts}

The $\sim 60$ sink formation events we observe over $15\kyr$ place
our simulation at the high end of the published range.
Early simulations with cosmological initial conditions found
2--5 fragments per halo within the first few hundred years
\citep{greif11, clark11}, while longer integrations over
$\sim 5\kyr$ recovered $\sim 10$ sinks \citep{stacy12}.
The radiation-hydrodynamic cosmological zoom-ins of
\citet{latif_khochfar22} across five minihaloes produced up to 23
stars per halo (55 total), with masses $0.1$--$37\,\Msun$ and a
slope $\mathrm{d}N/\mathrm{d}M \propto M^{-1.18}$.
The systematic study of \citet{wollenberg20} with 45 controlled
simulations (varying rotation and turbulence) found 5--20 sinks per
halo, confirming that fragmentation is highly stochastic and sensitive
to initial conditions.

Our fragment count likely reflects three compounding factors.
First, our $15\kyr$ integration is substantially longer than most
published runs, allowing many cycles of instability, collapse, and
re-accretion.
Second, cosmological initial conditions from the m12f progenitor
provide asymmetric, turbulent infall that continually replenishes
the outer disk and seeds new unstable structures.
Third, our mass-to-flux ratio remains comfortably supercritical
($\mu_\Phi \gg 1$) for most of the simulation, so magnetic fields
do not suppress fragmentation here as they do in some studies
\citep{stacy22, machida25}.
The high fragment count is not unprecedented: \citet{latif_khochfar22}
obtained comparable numbers with RHD (no MHD), and
\citet{prole22} showed that sink counts increase monotonically with
resolution without converging, suggesting our result may be a lower
bound.
\note{It would be worth quantifying: how many unique sink \textit{formation}
events occur vs the $\sim 30$ surviving at $15\kyr$? Stating the total
event count vs snapshot count separately would strengthen this comparison.}

% -----------------------------------------------------------
\subsection*{Role of the large-scale environment}

A key distinction between our simulation and most previous work is
the identity of the host halo.
At the epoch of first collapse ($z \sim 21$), the m12f progenitor is
a $\sim 10^5$--$10^6\,\Msun$ minihalo -- comparable in mass to those
studied in most other work -- but it is embedded in a large-scale
overdensity destined to grow to $M_\mathrm{vir} \sim 10^{12}\,\Msun$
by $z = 0$.
The most directly comparable simulation, \citet{meziani26}, targets
the m09 halo whose progenitor at $z \sim 14$ resides in a quieter
environment (reaching $\sim 2 \times 10^9\,\Msun$ at $z = 0$),
and finds a \textit{single} Pop~III protostar growing to
$\sim 27\,\Msun$ over $31\kyr$ -- in stark contrast to our
$\sim 60$-fragment cluster.
Although this comparison is a single pair of halos, the contrast
suggests that large-scale environment may be a key driver of whether
Pop~III star formation produces a single massive object or a
protostellar cluster.
The overdense filamentary infall geometry of the m12f progenitor
naturally seeds higher angular momentum, more complex disk structure,
and stronger turbulence than the more isolated m09 environment.
\citet{ho25} showed that infall into dark-matter potential wells
drives supersonic turbulence ($\mathcal{M} \sim 2$--$4$) across all
minihaloes; the m12f overdensity likely places our halo toward the
high-$\mathcal{M}$ end of this distribution.
Testing this environmental dependence statistically will require a
suite of cosmological RMHD zoom-ins spanning a range of $z=0$ halo
masses.
\note{Can we quantify the infall rate or Mach number in our simulation
and compare directly to Meziani+2026? A concrete number (e.g.\ higher
$\mathcal{M}$ or mass infall rate onto the disk) would make this
environmental argument much more compelling than a qualitative claim.}

% -----------------------------------------------------------
\subsection*{Stellar mass growth and accretion rates}

The total stellar mass of $\sim 200\,\Msun$ at $15\kyr$ is broadly
consistent with canonical Pop~III accretion rates.
\citet{stacy10} found $\dot{M} \sim 10^{-3}\,\Msun\,\mathrm{yr}^{-1}$
sustained over $\sim 5\kyr$; integrating this rate over $15\kyr$
gives $\sim 15\,\Msun$, far less than we find, reflecting our larger
sink count sharing a common reservoir rather than a single accreting
object.
Our time-averaged rate of $\sim 10^{-2}\,\Msun\,\mathrm{yr}^{-1}$
(across all sinks) is at the high end of the range
$10^{-4}$--$10^{-2}\,\Msun\,\mathrm{yr}^{-1}$ found by
\citet{hirano14} across 100 minihaloes, again consistent with a
particularly turbulent, high-infall-rate environment.

The power-law growth $M_* \propto t^{2.08}$ confirms the
$M_* \propto t^2$ scaling found by \citet{murray16} in simulations of
self-gravitating turbulent collapse, where the density inside the
sphere of influence converges to $\rho \propto r^{-3/2}$ and
Keplerian infall gives $\dot{M} \propto t$.
Our slightly super-quadratic exponent likely reflects the additional
contribution of continued cosmological infall sustaining the gas
reservoir at large radii, absent in the isolated box setup of
\citet{murray16}.
The POPSICLE simulations of \citet{sharda_menon25} produce a single
most-massive protostar of $\sim 65\,\Msun$ within $5\kyr$ in their
RMHD run; our total mass of $\sim 200\,\Msun$ is distributed across
$\sim 30$ surviving sinks at $15\kyr$, implying a characteristic
fragment mass of $\sim 5$--$7\,\Msun$, consistent with
fragmentation-induced starvation redistributing mass from the most
massive object to the cluster population \citep{prole22}.
\citet{sugimura20} and \citet{sugimura23} produced massive Pop~III
binaries with individual masses $>30\,\Msun$ at separations
$>2000\,\AU$; the two most massive sinks in our simulation -- which
dominate the mass budget throughout -- may represent an analogous
central binary, with the surrounding cluster population arising
from subsequent disk fragmentation.

% -----------------------------------------------------------
\subsection*{Star formation efficiency}

The star formation efficiency we find ($f_* \sim 50$--$80\%$ within
the disk at $15\kyr$) is high compared to both present-day values
($\sim 1$--$10\%$) and to other Pop~III simulations.
\citet{hirano_sakai23} report $\epsilon_\mathrm{III} \sim 10^{-2}$
in their large halo sample at the point of initial collapse; our
higher integrated value reflects the longer post-collapse integration
time and the lack of negative feedback mechanisms that would quench
accretion.
Radiative feedback is included in our simulation (M1 radiation
transport), but within the $15\kyr$ of evolution covered here no
protostar has yet grown sufficiently massive and contracted
sufficiently to launch an ionising H{\small II} region that would
reverse the accretion flow.
\citet{hosokawa11} found that UV feedback halts accretion at
$\sim 43\,\Msun$ for a single protostar on a $\sim 10^5$ yr
Kelvin-Helmholtz timescale; with our mass distributed across
many objects, no individual protostar yet dominates the radiative
output, and the global accretion efficiency remains high.
This suggests that the high SFE we observe is a transient feature of
the early accretion phase; longer integrations incorporating
protostellar contraction and UV feedback would be required to
determine final stellar masses and the termination of star formation.
\note{Confirm whether protostellar jets are active in the simulation
by $15\kyr$. If so, their momentum injection would be the dominant
early feedback channel before UV feedback, and should be mentioned
here. Meziani+2026 find jets launch when $\mu_\Phi \sim \mathrm{few}$,
which our inner disk approaches by $\sim 10\kyr$.}

% -----------------------------------------------------------
\subsection*{Protostellar mass function}

Our protostellar mass function (PMF) at $15\kyr$ has a high-mass
slope $\alpha = -0.99^{+0.14}_{-0.14}$, shallower than the Salpeter
slope ($-2.35$) but broadly consistent with other Pop~III simulations
that find top-heavy, log-flat distributions.
\citet{latif_khochfar22} find $\alpha \approx -1.18$ across five
halos with RHD; \citet{stacy16} and \citet{greif11} both report
roughly log-flat mass functions extending from sub-solar to
$\gtrsim 10\,\Msun$.
We emphasise that our PMF is a snapshot at $15\kyr$, not a final IMF:
all sinks are still actively accreting, no radiative feedback has
terminated accretion, and the population will be further shaped by
continued mergers, ejections, and eventual UV feedback
\citep{hosokawa11, hosokawa16, stacy16}.
\note{The two most massive sinks dominate the mass budget. What are
their individual masses at $15\kyr$? If either exceeds $\sim 20$--$30\,\Msun$
it would be useful to compare directly to Hosokawa+2011 (single-star
UV feedback limit) and Stacy+2016 (binary with feedback).
Also: are the most massive sinks still accreting rapidly, or has their
accretion stalled due to local gas depletion?}
The $\sim 52\%$ merger fraction and $\sim 23\%$ ejection fraction
we observe are consistent with the $\sim 50\%$ and $\sim 25\%$
values of \citet{stacy12} and the $\sim 30\%$ ejection rate of
\citet{wollenberg20}.
\citet{latif_khochfar22} report an unusually high ejection fraction
of $\sim 70\%$, which they attribute to the strong three-body
interactions driven by the radiative heating of the disk; the lower
fraction in our simulation may reflect the slower build-up of
radiation pressure in a multi-sink cluster where no single object
yet dominates.

% -----------------------------------------------------------
\subsection*{Magnetic field evolution}

Our mass-to-flux ratio evolves from $\mu_\Phi \sim 100$--$300$ at
early times to $\mu_\Phi \lesssim 1$ in the inner disk by $10\kyr$,
with magnetic energy growing by $\sim 4$ orders of magnitude.
This evolution is more dramatic than the $\mu_\Phi \sim 10$--$100$
reported by \citet{machida08} in idealised MHD simulations, likely
because our longer integration time allows more flux accumulation by
differential rotation winding and compression.
\citet{higashi24} demonstrated that the turbulent small-scale dynamo
amplifies fields to 30--50\% of equipartition during the primordial
collapse phase; our highly supercritical starting point and the
continued amplification we observe are consistent with this picture.

Despite this amplification, disk fragmentation in our simulation
remains vigorous throughout, in contrast to the strong suppression
found in some MHD simulations.
\citet{machida25} systematically varied initial field strength from
$B_0 = 10^{-20}$ to $10^{-4}$ G and found that even the weakest
field suppresses fragmentation and drives protostellar mergers within
the first $\sim 1000$~yr -- a timescale much shorter than our run.
\citet{stacy22} found only two secondary sinks in cosmological MHD
runs.
The key difference is that our simulation starts from a highly
supercritical configuration ($\mu_\Phi \sim 300$) and remains in the
fragmentation-permitting regime ($\mu_\Phi > 1$) for the majority of
the $15\kyr$.
Magnetic braking and flux-freezing act on a timescale long compared
to the free-fall time at the disk scale, so fragmentation proceeds
largely unimpeded while the field is amplified gradually from below.
Only toward the end of our simulation does the inner disk approach
the magnetically subcritical threshold; at this point,
\citet{meziani26} find that protostellar jets are launched, a phase
not yet reached in our $15\kyr$ window.
\note{It would be valuable to measure the magnetic energy power spectrum
at late times to assess whether turbulent dynamo action is contributing
to the amplification, or whether flux-freezing and winding alone account
for the $4$-decade magnetic energy growth. Our resolution may not be
sufficient to resolve the turbulent dynamo (see Higashi+2024 who find
$N_J \gtrsim 64$ cells per Jeans length is needed); explicitly stating
the Jeans resolution achieved would let the reader judge this.}
In contrast, \citet{sadanari21} and \citet{sadanari23} found that
ambipolar diffusion is dynamically subdominant in primordial gas
(diffusion heating always smaller than compressional heating), broadly
consistent with our non-ideal MHD treatment where ideal MHD dominates
throughout \citep{meziani26}.
\note{The tanaka24 citation used elsewhere in the paper (``turbulent
magnetic fields reduce but do not eliminate fragmentation'') should be
verified -- the exact paper identity is uncertain. If this is Sadanari
et al.\ 2024 (arXiv:2405.15045 on turbulent B-fields in first-star
disks), update the bib key and add the reference here.}
